% !PRINTED 89
\weberpage{89}
%<<CONT

7. \gesperrt{The number} $\varrho$ \gesperrt{of the invariants divisible by a prime} $p$
\gesperrt{must be equal to or less than the number} $\mu$
\gesperrt{of the index moduli, and it must be possible to form from the matrix}

\weq{24}{
\begin{gathered}
e_{1,1}\,i_{1,1}, \quad e_{2,1}\,i_{2,1}, \; \dots \; e_{\varrho,1}\,i_{\varrho,1} \\
e_{1,2}\,i_{1,2}, \quad e_{2,2}\,i_{2,2}, \; \dots \; e_{\varrho,2}\,i_{\varrho,2} \\
\cdot\ \cdot\ \cdot\ \cdot\ \cdot\ \cdot\ \cdot\ \cdot\ \cdot\ \cdot \\
e_{1,\mu}\,i_{1,\mu}, \quad e_{2,\mu}\,i_{2,\mu}, \; \dots \; e_{\varrho,\mu}\,i_{\varrho,\mu}
\end{gathered}
}

\gesperrt{at least one determinant of} $\varrho$ \gesperrt{rows that is not
divisible by} $p$\gesperrt{.}

This involves first of all a requirement on the $i_{k,h}$, that is,
in the last instance a condition on the number $m$.

All terms of any $\varrho$-rowed determinant of the matrix (24)
contain a factor of the form

\[
i_{1,h_1}\,i_{2,h_2}, \; \dots \; i_{\varrho,h_\varrho},
\]

in which $h_1, h_2, \dots h_\varrho$ denote any combination of $\varrho$
distinct numbers of the sequence $1, 2, \dots \mu$. Now if all
these products were divisible by $p$, then certainly all determinants of the
matrix (24) of $\varrho$ rows would also be divisible by $p$, and so there
must be at least one such combination that is not divisible by $p$.
But since the $i_{k,h}$ can only be powers of $p$, the equalities

\[
i_{1,h_1} = 1, \quad i_{2,h_2} = 1, \; \dots \; i_{\varrho,h_\varrho} = 1
\]

must hold. If we now go back to the meaning of the $i_{k,h}$ in (11), we
can express this last requirement on the degree $m$ as
follows:

f) \gesperrt{If a prime} $p$ \gesperrt{divides} $\varrho$
\gesperrt{invariants} $i_1, i_2, \dots i_\varrho$\gesperrt{, then the number}
$\mu$ \gesperrt{of the index moduli must be equal to or greater than}
$\varrho$\gesperrt{, and it must be possible to select} $\varrho$ \gesperrt{of these
index moduli} $c_{h_1}, c_{h_2}, \dots c_{h_\varrho}$ \gesperrt{in such a way
that} $c_{h_1}$ \gesperrt{is divisible by} $i_1$, $c_{h_2}$
\gesperrt{by} $i_2 \dots$, $c_{h_\varrho}$ \gesperrt{by} $i_\varrho$ $^1$).

\noindent\rule{4cm}{0.4pt}

\footnotesize
$^1$) The proof of the fact that this requirement can always be
satisfied in infinitely many ways rests on the theorem
of number theory that among the numbers of an arithmetic progression
whose initial term and difference are without common divisor there occur
infinitely many primes -- a theorem which up to now has been established only by
\normalsize
%CONT>>
